\section{Fatigue damage and the Fatigue Damage Spectrum}\label{sec:damage}

\subsection{Basquin's law and Miner's rule}

A material loaded by cycles of constant stress amplitude $S$ fails after $N(S)$ cycles. In the range
of interest the S--N curve is a straight line in logarithmic scales (Basquin's law):
\begin{equation}
  N(S)\,S^b = C ,
\end{equation}
with the exponent $b$ (typically 3 to 12; 8 is a common default for structural details) and the
constant $C$. A cycle of amplitude $S$ consumes the fraction $1/N(S) = S^b/C$ of the life.
Miner's rule adds these fractions: a history containing $n_c$ cycles of amplitude $S_c$ has the damage
\begin{equation}\label{eq:miner}
  D = \sum_c \frac{n_c}{N(S_c)} = \frac1C\sum_c n_c S_c^b ,
\end{equation}
and failure is expected when $D$ reaches 1. With $S = Kz$ (Section~\ref{sec:oscillator}), the damage
of an oscillator is a sum over the cycles of its relative displacement:
\begin{equation}\label{eq:minerz}
  D = \frac{K^b}{C}\sum_c n_c\, A_c^b ,\qquad A_c = \text{amplitude (half range) of cycle } c .
\end{equation}
Only the ratio of two damages matters in test tailoring, so $K$ and $C$ are set to 1 by default.

\subsection{Rainflow counting}\label{sec:rainflow}

A random response does not consist of clean cycles; \emph{rainflow counting} decomposes it into
cycles in the way that is consistent with the stress--strain hysteresis of the material. The program
proceeds in two steps.
\begin{enumerate}[nosep]
\item \emph{Turning points.} Only the local maxima and minima of the sampled response matter; the
samples in between and plateaus are removed, leaving an alternating sequence $s_1, s_2, \dots$.
\item \emph{Four-point rule.} Take four consecutive turning points $s_{l}, s_{l+1}, s_{l+2},
s_{l+3}$ and the three ranges between them. If the inner range $|s_{l+2}-s_{l+1}|$ is not larger than
either outer range, the inner pair forms a closed cycle (a small loop inside a larger one): it is
counted with amplitude $|s_{l+2}-s_{l+1}|/2$ and the two points are removed. The rule is applied until
no more cycles close. The remaining sequence (the \emph{residue}) is counted as half cycles, one per
range, with weight $n_c = \tfrac12$.
\end{enumerate}
The program applies the four-point rule to all disjoint candidate pairs of the sequence at once, which
is fast in NumPy, and finishes long monotone decays with a sequential stack. Its damage sums are equal
to those of the sequential three-point algorithm of the ASTM E1049 standard (checked to $10^{-10}$ in
the tests).

\subsection{Fatigue Damage Spectrum and Maximum Response Spectrum}\label{sec:fds}

\begin{result}[FDS and MRS]
For a base acceleration $x$ of duration $T$, the Fatigue Damage Spectrum is the damage~\eqref{eq:minerz}
of the rainflow cycles of the response $z_{\fn}$ of each oscillator of the grid,
\begin{equation}\label{eq:fds}
  D(\fn) = r\,\frac{K^b}{C}\sum_{c\,\in\,\mathrm{rainflow}(z_{\fn})} n_c A_c^b ,
\end{equation}
and the Maximum Response Spectrum is the largest pseudo-acceleration,
\begin{equation}
  \mathrm{MRS}(\fn) = \wn^2\max_t|z_{\fn}(t)|.
\end{equation}
\end{result}
The factor $r$ is the number of times the signal is played: a record that represents a life
$\Tl$ longer than its duration $T$ has $r = \Tl/T$ (rounded to an integer when $\Tl\ge T$, i.e.\ the
record is considered repeated), and a test profile that is shorter than the test duration is played $r$
times (Section~\ref{sec:practice}). Damage adds over repetitions; the maximum does not, so the MRS has
no factor $r$. The MRS is a check of the peak stresses: a test whose MRS exceeds that of the life by
much may fail the component by overload rather than by fatigue.

\subsection{Rice's formulas for a Gaussian response}\label{sec:rice}

For a stationary Gaussian response the expected damage can be computed from the spectral moments.

\emph{Zero up-crossing rate.} The response crosses zero upwards in the interval $(t,t+\dd t)$ if
$z(t)\in(-\dot z\,\dd t, 0)$ with $\dot z>0$. The probability of this event is $\int_0^\infty
p(0,\dot z)\,\dot z\,\dd t\,\dd\dot z$, where $p(z,\dot z)$ is the joint density of $z$ and $\dot z$.
For a stationary process $z$ and $\dot z$ are uncorrelated, because $\E[z\dot z] =
\tfrac12\,\dd\E[z^2]/\dd t = 0$, and being jointly Gaussian they are independent:
$p(0,\dot z) = \frac{1}{\sqrt{2\pi}\sigma_z}\cdot\frac{1}{\sqrt{2\pi}\sigma_{\dot z}}
e^{-\dot z^2/(2\sigma_{\dot z}^2)}$. Dividing the probability by $\dd t$ and using
$\int_0^\infty \dot z\,e^{-\dot z^2/(2\sigma_{\dot z}^2)}\dd\dot z = \sigma_{\dot z}^2$,
\begin{equation}\label{eq:nu0}
  \nu_0^+ = \frac{\sigma_{\dot z}^2}{2\pi\sigma_z\sigma_{\dot z}} = \frac{\sigma_{\dot z}}{2\pi\sigma_z}
  = \sqrt{\frac{m_2}{m_0}} ,
\end{equation}
using $\sigma_{\dot z} = 2\pi\sqrt{m_2}$ and $\sigma_z = \sqrt{m_0}$ (Section~\ref{sec:respvar}).
The same reasoning applied to $\dot z$ (a maximum of $z$ is a down-crossing of zero by $\dot z$) gives
the rate of peaks $\nu_p = \sqrt{m_4/m_2}$. The ratio $\gamma = \nu_0^+/\nu_p = m_2/\sqrt{m_0m_4}$ is
the \emph{irregularity factor}: $\gamma=1$ means one peak per zero crossing (narrow band).

\emph{Amplitudes of a narrow-band response.}\label{sec:rayleigh} The response of a lightly damped
oscillator is narrow band: a sinusoid at $\fn$ whose amplitude changes slowly,
$z = A\cos(\wn t+\psi)$. Then $\dot z/\wn \approx -A\sin(\wn t+\psi)$ and the amplitude is
$A = \sqrt{z^2+(\dot z/\wn)^2}$. At a given time $z$ and $\dot z/\wn$ are independent zero-mean
Gaussian variables with the same variance $\sigma^2 = \sigma_z^2$ (equipartition,
Section~\ref{sec:envelope}). In polar coordinates $(A,\psi)$ their joint density
$\frac{1}{2\pi\sigma^2}e^{-A^2/(2\sigma^2)}$ times the Jacobian $A$, integrated over $\psi$, gives the
\emph{Rayleigh} density of the amplitude:
\begin{equation}
  p(A) = \frac{A}{\sigma^2}\,e^{-A^2/(2\sigma^2)},\qquad A\ge0 .
\end{equation}
Its moment of order $b$ follows with the substitution $u = A^2/(2\sigma^2)$, $A = \sqrt{2\sigma^2u}$,
$\dd u = A\,\dd A/\sigma^2$:
\begin{equation}\label{eq:rayleighmoment}
  \E[A^b] = \int_0^\infty A^b\,\frac{A}{\sigma^2}e^{-A^2/(2\sigma^2)}\,\dd A
  = \int_0^\infty (2\sigma^2u)^{b/2}e^{-u}\,\dd u = (2\sigma^2)^{b/2}\,\Gamma\!\left(1+\tfrac b2\right).
\end{equation}

\emph{Expected damage.} A narrow-band response has one cycle per period, $\nu_0^+T \approx \fn T$
cycles in the duration $T$, each with a Rayleigh amplitude. By~\eqref{eq:minerz} the expected damage is
the number of cycles times $\E[A^b]$:
\begin{result}[Gaussian narrow-band damage (Rice--Lalanne)]
\begin{equation}\label{eq:dgauss}
  D_G(\fn) = \frac{K^b}{C}\,\fn\,T\,\bigl(2\sigma_z^2\bigr)^{b/2}\,\Gamma\!\left(1+\tfrac b2\right),
  \qquad \sigma_z^2 = \sum_j W_{ij}G_j .
\end{equation}
\end{result}
Solving for the variance gives the response variance that produces a prescribed damage $D$ in the
duration $T$:
\begin{equation}\label{eq:targetvar}
  \sigma_z^2 = \frac12\left(\frac{D\,C}{K^b\,\fn\,T\,\Gamma(1+b/2)}\right)^{2/b}.
\end{equation}
For wide-band responses (for instance the quasi-static response of a high-frequency oscillator to a
strong low-frequency input) the cycles are not one per period and their amplitudes are not Rayleigh
distributed. Dirlik's empirical distribution of the rainflow ranges (Appendix~\ref{app:dirlik})
covers this case; it reduces to~\eqref{eq:dgauss} in the narrow-band limit $\gamma\to1$. MMVSD uses
it (Section~\ref{sec:damagemodel}).

\subsection{Rainflow on samples: peaks between samples}\label{sec:sampling}

Rainflow counting works on samples. The true maximum of a cycle generally falls between two samples;
the largest sample of the half cycle is the one nearest to it. For a response of frequency $f$ the
nearest sample lies within $\pm\Delta t/2$ of the true peak, at a phase offset $\vartheta = 2\pi f
\delta$ uniformly distributed in $(-\pi f/f_s, \pi f/f_s)$, and the counted peak is $A\cos\vartheta$
instead of $A$. A range is measured between a maximum and the following minimum, with independent
offsets $\vartheta_1,\vartheta_2$; its amplitude is $A(\cos\vartheta_1+\cos\vartheta_2)/2$, and its
damage is reduced on average by the factor
\begin{result}[Sampling factor]
\begin{equation}\label{eq:sampling}
  c_b(\rho) = \E\!\left[\left(\frac{\cos\vartheta_1+\cos\vartheta_2}{2}\right)^{\!b}\right]
  = \frac{1}{(\pi\rho)^2}\int_0^{\pi\rho}\!\!\int_0^{\pi\rho}
  \left(\frac{\cos\vartheta_1+\cos\vartheta_2}{2}\right)^{\!b}\dd\vartheta_1\,\dd\vartheta_2 ,
  \qquad \rho = \frac{f}{f_s},
\end{equation}
\end{result}
(the integrand is even in each angle, so the average over $(-\pi\rho,\pi\rho)$ equals the average
over $(0,\pi\rho)$). For $b=8$: $c_b = 0.98$ at $\rho = 0.04$, $0.86$ at $\rho = 0.11$ and $0.54$ at
$\rho=0.23$. The program evaluates~\eqref{eq:sampling} by Gauss--Legendre quadrature on a table of
$\rho$. Together with the digital response of Section~\ref{sec:digital} this explains why a Gaussian
design underestimates the measured damage loss at high natural frequencies when $f_s$ is close to
$2.5f_{\max}$; MMVSD includes both effects in its model.

\subsection{Why the kurtosis of the response matters}\label{sec:kzdamage}

Let the response be Gaussian given a slowly varying local variance $s(t)$ (Section~\ref{sec:modulated}).
Over a stretch of time where $s$ is constant, the damage rate is the narrow-band
rate~\eqref{eq:dgauss} with $\sigma_z^2 = s$. Averaging over time,
\begin{equation}
  D = \frac{K^b}{C}\fn T\,\Gamma\!\left(1+\tfrac b2\right) 2^{b/2}\,\tavg{s^{b/2}}
  = D_G\bigl(\tavg{s}\bigr)\;\frac{\tavg{s^{b/2}}}{\tavg{s}^{b/2}} ,
\end{equation}
where $D_G(\tavg s)$ is the Gaussian damage of a stationary response with the same variance.
The factor $\tavg{s^{b/2}}/\tavg{s}^{b/2}$ is at least 1 (Jensen's inequality, since $s^{b/2}$ is
convex for $b\ge2$). For $b=4$ it equals $\tavg{s^2}/\tavg{s}^2 = \kappa_z/3$ by~\eqref{eq:kmod}; for
larger $b$ it depends on higher moments of $s$ but grows with $\kappa_z$. Two signals with the same PSD
can therefore produce very different damage if their responses have different kurtosis, and a test
that reproduces the FDS with a different $\kappa_z(\fn)$ reproduces it with a different mix of cycles
(many moderate cycles instead of a few large ones). This is the reason for prescribing $\kappa_z$
together with the FDS (Part~III).
